% !TEX program = xelatex
\documentclass[a4paper, 11pt]{ctexart}
\usepackage[margin=2cm]{geometry}
\usepackage{xcolor}
\usepackage{booktabs}
\usepackage{titlesec}
\usepackage{tcolorbox}
\usepackage{enumitem}
\usepackage{hyperref}
\usepackage{float}
\usepackage{array}
\usepackage{amsmath}
\usepackage{amssymb}
\usepackage{pifont}

\definecolor{tsinghua}{RGB}{102, 8, 116}
\definecolor{bglight}{RGB}{247, 245, 248}
\definecolor{okgreen}{RGB}{34, 139, 34}
\definecolor{warnred}{RGB}{180, 30, 30}
\titleformat{\section}[block]{\Large\bfseries\color{tsinghua}}{\thesection}{1em}{}
\titleformat{\subsection}[block]{\large\bfseries\color{tsinghua!80!black}}{\thesubsection}{1em}{}
\linespread{1.35}
\setlist[itemize]{itemsep=4pt, parsep=0pt, topsep=4pt}
\hypersetup{colorlinks=true, linkcolor=tsinghua, urlcolor=tsinghua}

\title{\vspace{-1.5cm}\color{tsinghua}\textbf{方向 C：工业故障诊断中的检索增强生成\\——新框架、新指标与首个中文评测基准}\\\Large{—— ``鸿雁计划''科创基金细化企划书 v3}}
\author{\textbf{汇报人：纪文龙} \\ \small 清华大学行健书院 2024级本科生}
\date{\small 2026年4月}

\begin{document}
\maketitle
\thispagestyle{empty}

\begin{tcolorbox}[colback=tsinghua!5, colframe=tsinghua!50, title=\textbf{一句话概括}, fonttitle=\bfseries]
针对通用RAG在工业文档检索中的\textbf{三类特殊失败模式}（因果链断裂、跨文档推理失败、术语不对等），提出\textbf{知识图谱增强的三路混合检索框架}与\textbf{两个原创评估指标}（Chain Completeness、Industrial Faithfulness Score），并构建\textbf{首个中文动力设备故障诊断问答评测数据集}。通过三组严格对照实验量化各项改进幅度，产出1篇SCI论文（目标Q1）。
\end{tcolorbox}

% ================================================================
\section{核心定位：RAG是工具，研究问题才是目标}
% ================================================================

\subsection{项目的基本逻辑}

在目前的开源社区支持下（LangChain、LlamaIndex、Chroma等），搭建一个基础的RAG系统已经\textbf{不难、不慢}——这不是我们的研究贡献。

我们的研究贡献在于：在搭好RAG之后，围绕动力设备故障诊断这一场景，\textbf{回答三个有科学价值的研究问题}，并用严格的对照实验给出\textbf{可量化的效果}。

\begin{tcolorbox}[colback=bglight, colframe=tsinghua!30, arc=2mm]
\textbf{类比}：PyTorch是工具，不是论文。但用PyTorch训练模型、发现规律、解决新问题——那是论文。同理，RAG是工具，用RAG回答领域内尚未解决的问题——那才是论文。

\medskip
\textbf{本研究的论文故事线}：\\
\centerline{识别问题（通用RAG在工业领域的失败模式） → 提出方案（新框架+新指标） → 实验验证（三组对照实验） → 贡献资源（首个中文Benchmark）}
\end{tcolorbox}

\subsection{为什么这个方向``对将来帮助更大''？}

\begin{itemize}
    \item \textbf{技术栈的未来价值}：LLM/RAG/Agent 是未来 5--10 年工业AI的核心技术路线。掌握它比学一种具体的CNN架构更有长期回报。
    \item \textbf{技能的可迁移性}：学会构建领域知识库、设计AI问答评估体系，这些能力横跨所有行业——医疗、法律、制造、能源都需要。
    \item \textbf{个人技术优势}：Web开发能力在系统演示和原型搭建上有直接用武之地。
\end{itemize}


% ================================================================
\section{当前技术前沿与文献空白}
% ================================================================

\subsection{2024--2025年，这个方向发生了什么？}

\begin{table}[H]
\centering
\renewcommand{\arraystretch}{1.2}
\small
\begin{tabular}{p{3.5cm} p{5cm} p{4.5cm}}
    \toprule
    \textbf{趋势} & \textbf{代表性工作} & \textbf{对我们的启示} \\
    \midrule
    知识图谱增强RAG & KBS 2025：Graph RAG用于高铁转向架故障诊断（DOI: 10.1016/j.knosys.2025.114855） & 证明KG+RAG+故障诊断可发Q1 \\
    多Agent诊断系统 & MAKG（2026）：多Agent+KG协同，90.1\%推理准确率 & 多Agent路线已有先例 \\
    领域检索评估 & PARAM（arXiv 2025）：振动数据序列化+RAG自动维修建议 & 聚焦评估和可信度是差异化方向 \\
    制造业GraphRAG & Electronics 2025：Document GraphRAG用于制造业QA & Q2级别难度可控 \\
    \bottomrule
\end{tabular}
\end{table}

\subsection{我们能填补什么空白？}

\begin{tcolorbox}[colback=bglight, colframe=tsinghua!30, arc=2mm]
上述已发表工作有三个共同缺陷：\textbf{（1）没有针对工业文档检索的失败案例分析}——都是直接报好结果，不解释``为什么通用RAG在工业领域不够好''；\textbf{（2）没有多步Agent与单轮RAG的严格对照实验}——无法量化Agent到底带来多大收益；\textbf{（3）没有工业安全场景的专用评估指标}——都用通用NLP指标，不考虑安全约束。

此外，\textbf{中文动力设备故障诊断的问答评测数据集目前为空白}。
\end{tcolorbox}


\subsection{与已发表直接可比论文的差异化对比}

下表从五个关键维度，逐一说明本研究相比三篇已发表/已接收的直接可比论文\textbf{具体多做了什么}：

\begin{table}[H]
\centering
\renewcommand{\arraystretch}{1.25}
\small
\begin{tabular}{p{2.8cm} p{2.4cm} p{2.4cm} p{2.4cm} p{2.8cm}}
    \toprule
    \textbf{对比维度} & \textbf{KBS 2025} \newline 高铁转向架 & \textbf{Electronics 2025} \newline 制造业文档QA & \textbf{MAKG 2026} \newline 工业维护 & \textbf{本研究} \\
    \midrule
    检索策略 & KG增强 & KG增强 & 多Agent分工 & \textcolor{okgreen}{\textbf{KG+BM25+向量三路混合+Reranker}} \\
    多步推理Agent & {\color{warnred}\ding{55}} & {\color{warnred}\ding{55}} & \checkmark（无对照） & \textcolor{okgreen}{\textbf{\checkmark + 三级难度对照实验}} \\
    评估指标 & 通用NLP指标 & 通用NLP指标 & 通用NLP指标 & \textcolor{okgreen}{\textbf{IFS + Chain Completeness（原创）}} \\
    中文Benchmark & {\color{warnred}\ding{55}} & {\color{warnred}\ding{55}} & {\color{warnred}\ding{55}} & \textcolor{okgreen}{\textbf{\checkmark \, 首个（200+问答对）}} \\
    安全约束考量 & {\color{warnred}\ding{55}} & {\color{warnred}\ding{55}} & {\color{warnred}\ding{55}} & \textcolor{okgreen}{\textbf{\checkmark \, 安全拒答率}} \\
    \bottomrule
\end{tabular}
\caption{与三篇直接可比论文的系统性对比。KBS 2025 = Knowledge-Based Systems, Vol.331; Electronics 2025 = Electronics, Vol.14(11); MAKG = 多Agent知识图谱框架（2026预印本）。}
\end{table}

\begin{tcolorbox}[colback=tsinghua!5, colframe=tsinghua!50, arc=2mm]
\textbf{关键结论}：与发在KBS（Q1, IF$\approx$8.1）上的最新同类工作相比，本研究在\textbf{多步推理对照实验}、\textbf{原创评估指标}、\textbf{中文Benchmark}、\textbf{安全约束}四个维度均有明确的增量贡献。若一篇只做了``KG增强检索''的工作可以发Q1，我们的方案发Q1是完全合理的预期。
\end{tcolorbox}


% ================================================================
\section{三个研究问题}
% ================================================================

\subsection{RQ1：工业文档检索有什么特殊困难？针对性改进能提升多少？}

\textbf{研究动机}：我们不会凭空说``需要知识图谱增强''。而是先搭一个基础RAG，观察它在动力设备场景中\textbf{具体在哪些类型的问题上表现不好}，分析原因，再针对性地提出改进方案。这符合科学研究的基本逻辑：\textbf{观察 → 假设 → 实验 → 结论}。

\textbf{预期发现的三类典型检索失败}：
\begin{itemize}
    \item \textbf{因果链断裂}：用户问``汽轮机振动大怎么排查？''，基础RAG找到了振动相关段落，但没能把``不平衡 → 对中检查 → 动平衡''这条完整的排查链路检索出来。原因：因果关系分散在多个文档段落中。
    \item \textbf{跨文档推理失败}：答案需要综合3份不同文档的信息（设备手册+历史案例+行业标准），但基础RAG只检索到了其中一份。
    \item \textbf{术语不对等}：用户说``机组抢振''，文档中写的是``转子不平衡导致的强制振动''——语义相同但措辞不同，向量检索匹配不上。
\end{itemize}

\textbf{针对性改进方案}：
\begin{enumerate}
    \item 针对因果链断裂 → 构建\textbf{故障因果知识图谱}（设备→部件→故障→原因→方案），用图遍历补全因果链
    \item 针对跨文档推理 → \textbf{三路混合检索}：向量检索 + BM25关键词检索 + 知识图谱子图查询，结果加权融合后用Reranker精排
    \item 针对术语不对等 → 构建\textbf{领域同义词表} + LLM驱动的\textbf{Query改写}
\end{enumerate}

\textbf{实验设计}：
\begin{itemize}
    \item 对照组：基础RAG（纯向量检索 + 固定分块）
    \item 实验组：改进RAG（知识图谱 + 混合检索 + Query改写）
    \item 测试集：100--200道领域问题（事实型/诊断型/方法型三类）
    \item 指标：\textbf{Hit Rate@5}、\textbf{MRR}、以及我们提出的原创指标\textbf{Chain Completeness}
\end{itemize}

\begin{tcolorbox}[colback=tsinghua!5, colframe=tsinghua!40, title=\textbf{原创指标：Chain Completeness (CC)}, fonttitle=\bfseries]
\textbf{动机}：传统检索评估指标（Hit Rate、MRR）只衡量``答案文档有没有被找到''，但不关心检索结果是否覆盖了\textbf{完整的诊断推理链}。在故障诊断中，仅找到``振动大''的文档是不够的，还需要找到``不平衡→对中→动平衡''的完整排查路径。现有指标无法衡量这一维度。

\textbf{定义}：给定一个诊断问题$q$，其标准答案包含一条因果链$S_{\text{gold}}(q) = \{s_1, s_2, \ldots, s_n\}$（例如``症状→可能原因→排查步骤→修复方案''各环节）。检索系统返回文档集$D_{\text{ret}}(q)$。Chain Completeness定义为：
\[
CC(q) = \frac{\bigl|\{s_i \in S_{\text{gold}}(q) \mid \exists\, d_j \in D_{\text{ret}}(q),\; s_i \sqsubseteq d_j \}\bigr|}{|S_{\text{gold}}(q)|}
\]
其中$s_i \sqsubseteq d_j$表示因果链步骤$s_i$的语义被文档$d_j$覆盖（通过NLI模型或人工判断确定）。

\textbf{系统级指标}：$\overline{CC} = \frac{1}{|Q|}\sum_{q \in Q} CC(q)$，即测试集上所有问题的平均Chain Completeness。

\textbf{为什么是创新贡献}：这是首个专门衡量``检索结果是否覆盖完整诊断推理链''的评估指标，\textbf{填补了RAG评估在工业故障诊断场景中的空白}。
\end{tcolorbox}

\textbf{可衡量的效果}：改进后检索准确率比基础RAG提升 $\geq 5\%$（参考KBS 2025论文，类似改进通常在10--20\%量级）。Chain Completeness的提升幅度预期更大，因为基础RAG在因果链覆盖上存在系统性缺陷。


\subsection{RQ2：多步Agent推理比单轮问答好多少？}

\textbf{研究动机}：真实的故障排查本质上是\textbf{多步信息收集 + 假设检验 + 逐步排除}的过程。一线工程师请教老师傅时，老师傅会先问``什么时候出现的？''''负荷多少？''——这种\textbf{迭代式追问}是单轮RAG无法实现的。我们的核心问题不是``能不能搭Agent''（工程问题），而是``Agent比不用Agent\textbf{好多少}''（研究问题）。

\textbf{Agent工具集设计}：

\begin{table}[H]
\centering
\renewcommand{\arraystretch}{1.2}
\small
\begin{tabular}{p{3.5cm} p{4cm} p{5cm}}
    \toprule
    \textbf{工具名} & \textbf{功能} & \textbf{实现方式} \\
    \midrule
    search\_knowledge & 知识库检索 & RQ1的混合检索 \\
    get\_fault\_chain & 获取故障因果链 & 知识图谱子图查询 \\
    query\_specs & 查设备参数 & 结构化数据库SQL查询 \\
    ask\_user & 追问用户细节 & LLM生成追问问题 \\
    generate\_diagnosis & 综合诊断 & LLM推理+引用来源 \\
    \bottomrule
\end{tabular}
\end{table}

\textbf{实验设计——三级难度 $\times$ 三种系统的对照实验}：

\begin{table}[H]
\centering
\renewcommand{\arraystretch}{1.2}
\small
\begin{tabular}{p{1.5cm} p{4cm} p{3.5cm} p{3.5cm}}
    \toprule
    \textbf{难度} & \textbf{问题特征} & \textbf{示例} & \textbf{预期结果} \\
    \midrule
    简单 & 单因素，答案在单篇文档中 & ``XX型轴承额定温升?'' & 三种系统差异小 \\
    中等 & 需结合2--3个信息源 & ``3号机组振动偏高的可能原因?'' & Agent小幅领先 \\
    复杂 & 多因素交叉、信息不完整 & ``异常声响+排气温度偏高，如何排查?'' & \textbf{Agent显著领先} \\
    \bottomrule
\end{tabular}
\end{table}

\begin{itemize}
    \item \textbf{基线A}：纯大模型直接回答（不检索）
    \item \textbf{基线B}：单轮RAG（检索一次 → 生成答案）
    \item \textbf{实验组C}：多步Agent（可检索、追问、调用工具、多轮推理）
\end{itemize}

\textbf{评估方式}：
\begin{itemize}
    \item 诊断准确率（与标准答案对比）
    \item 推理合理性（邀请刘老师/课题组博士生盲评打分，1--5分量表）
    \item 按难度分层统计——\textbf{核心预期结论}：``在复杂多因素故障场景中，多步Agent的诊断准确率比单轮RAG提升 $Y\%$''
    \item 计算评分者间一致性（Cohen's Kappa），确保专家评估的可靠性
\end{itemize}



\subsection{RQ3：工业安全场景需要什么样的评估指标？}

\textbf{研究动机}：现有RAG评估框架（RAGAS等）只关注``答案和文档是否一致''，不关注\textbf{工业安全约束}。但在动力设备领域，一个``看似合理但参数错误''的回答可能导致安全事故。通用指标无法捕捉这类风险。

\textbf{具体研究步骤}：

\begin{enumerate}
    \item \textbf{发现问题}：用通用RAGAS指标（Faithfulness、Answer Relevancy）评估系统，同时请领域专家评分。对比两者结果，\textbf{找出``RAGAS给高分但专家说不对''的案例}——这些案例揭示了通用指标在工业领域的盲点。
    
    \item \textbf{提出新指标——Industrial Faithfulness Score (IFS)}：
    
    \item \textbf{验证指标有效性}：以专家评分为ground truth，计算IFS与专家评分的\textbf{Spearman相关系数}，并与RAGAS的通用指标对比。
\end{enumerate}

\begin{tcolorbox}[colback=tsinghua!5, colframe=tsinghua!40, title=\textbf{原创指标：Industrial Faithfulness Score (IFS)}, fonttitle=\bfseries]
\textbf{动机}：通用忠实度指标（如RAGAS Faithfulness）仅检查``回答是否基于检索文档''，但不检查\textbf{参数是否正确}、\textbf{因果逻辑是否成立}、\textbf{超出知识范围时是否安全拒答}。在安全攸关的工业场景中，这三类错误的后果远比一般NLP任务严重。

\textbf{定义}：给定系统对问题$q$的回答$a$，Industrial Faithfulness Score定义为三个子维度的加权组合：
\[
\text{IFS}(a, q) = \alpha \cdot \text{ParamAcc}(a) + \beta \cdot \text{CausalAlign}(a) + \gamma \cdot \text{SafeRef}(a, q)
\]

三个子维度分别为：

\begin{itemize}
    \item \textbf{参数一致性 $\text{ParamAcc}(a)$}：回答中引用的数值参数（温度、转速、压力等）与知识库记录的匹配比例。
    \[
    \text{ParamAcc}(a) = \frac{|\{\text{回答中的参数值}\} \cap \{\text{知识库中的正确值}\}|}{|\{\text{回答中的参数值}\}|}
    \]
    
    \item \textbf{因果逻辑一致性 $\text{CausalAlign}(a)$}：诊断链条中的``原因→结果''关系是否与知识图谱中的因果边一致。
    \[
    \text{CausalAlign}(a) = \frac{|\{\text{回答中的因果对}\} \cap \{\text{图谱中的因果边}\}|}{|\{\text{回答中的因果对}\}|}
    \]
    
    \item \textbf{安全拒答 $\text{SafeRef}(a, q)$}：面对超出知识范围的问题，系统正确拒绝回答（而非编造）的指示函数。
    \[
    \text{SafeRef}(a, q) = \begin{cases} 1, & \text{若}\ q \notin \text{知识范围}\ \text{且系统正确拒答} \\ 1, & \text{若}\ q \in \text{知识范围}\ \text{且系统正常作答} \\ 0, & \text{其他（编造或遗漏）} \end{cases}
    \]
\end{itemize}

权重$\alpha, \beta, \gamma$通过与专家评分的多元回归分析确定，使IFS与专家判断的相关性最大化。

\textbf{为什么是创新贡献}：这是首个面向\textbf{安全攸关工业场景}的RAG忠实度量化指标，将通用的``是否基于文档''扩展为``参数是否正确、逻辑是否成立、不确定时是否安全拒答''三个维度，\textbf{填补了工业垂直领域幻觉检测的空白}。
\end{tcolorbox}

\textbf{可衡量的效果}：预期结论——``IFS与专家评分的Spearman相关系数为 $r_s \approx 0.85$，而通用RAGAS Faithfulness指标仅为 $r_s \approx 0.60$''。


% ================================================================
\section{论文贡献总结}
% ================================================================

\begin{table}[H]
\centering
\renewcommand{\arraystretch}{1.3}
\small
\begin{tabular}{p{0.5cm} p{4.5cm} p{2cm} p{1.5cm} p{4cm}}
    \toprule
    \textbf{\#} & \textbf{贡献内容} & \textbf{贡献类型} & \textbf{创新度} & \textbf{已有工作未覆盖} \\
    \midrule
    1 & 系统分析工业文档检索的三类特殊困难，提出图谱增强三路混合检索策略并验证有效性 & 方法 + 实验 & 中--高 & KBS 2025仅用KG增强，未做失败案例分析和三路融合 \\
    2 & 提出\textbf{Chain Completeness}指标：首个衡量检索结果对诊断因果链覆盖度的评估指标 & 指标/方法论 & 高 & 所有已有工作均使用通用检索指标 \\
    3 & 多步Agent vs 单轮RAG的\textbf{三级难度对照实验}，量化Agent在复杂场景中的改进幅度 & 实验 & 中 & MAKG有多Agent但\textbf{无对照实验} \\
    4 & 提出\textbf{IFS指标}：首个面向安全攸关工业场景的RAG忠实度量化方法 & 指标/方法论 & 高 & 所有已有工作均无安全约束评估 \\
    5 & 构建首个\textbf{中文动力设备故障诊断问答评测数据集}（200+问答对，三种难度$\times$三种类型） & 数据集/资源 & 很高 & \textbf{目前完全空白} \\
    \bottomrule
\end{tabular}
\end{table}

\begin{tcolorbox}[colback=bglight, colframe=tsinghua!30, arc=2mm]
\textbf{发表目标}：投向 \textit{Expert Systems with Applications}（IF$\sim$8.5）或 \textit{Knowledge-Based Systems}（IF$\sim$8.1）。这两个Q1期刊在2024--2025年大量接收KG+RAG+工业应用类论文（如KBS 2025已发表``Graph RAG for Train Bogies''同类工作），\textbf{明确接受系统论文和数据集贡献论文}。

\textbf{双轨策略}：
\begin{itemize}[itemsep=2pt]
    \item \textbf{冲击Q1（KBS / ESWA）}：五项贡献齐备，重点突出IFS新指标和Benchmark数据集
    \item \textbf{稳拿Q2（Applied Sciences / Electronics）}：即使RQ2/RQ3做得不够深，仅凭贡献1（检索改进实验）+ 贡献5（中文Benchmark）已足够
\end{itemize}
\end{tcolorbox}


% ================================================================
\section{预期成果}
% ================================================================

\begin{itemize}
    \item \textbf{1篇SCI论文}：冲击Q1（KBS/ESWA），保底Q2（Applied Sciences/Electronics）。
    \item \textbf{2个原创评估指标}：Chain Completeness（检索因果链覆盖度）和IFS（工业忠实度评分），可供后续研究者在工业RAG领域复用。
    \item \textbf{1个评测数据集}：首个中文动力设备故障诊断QA Benchmark（200+问答对），按学术规范标准化后开源，预期成为本研究\textbf{最具长期引用价值}的产出。
    \item \textbf{1套演示系统}：Web界面的对话式诊断助手——输入故障描述，输出带引用来源的诊断建议 + Agent推理过程可视化。
    \item \textbf{鸿雁计划结题}：完整结题报告 + 上述成果。
\end{itemize}


% ================================================================
\section{风险评估与缓解策略}
% ================================================================

\begin{table}[H]
\centering
\renewcommand{\arraystretch}{1.3}
\small
\begin{tabular}{p{4.5cm} p{1.5cm} p{7cm}}
    \toprule
    \textbf{风险} & \textbf{概率} & \textbf{缓解策略} \\
    \midrule
    审稿人质疑``搭了个系统，没有方法创新'' & 中 & 靠两个原创指标（CC和IFS）和首个中文Benchmark对冲——数据集贡献和指标贡献是审稿人无法否认的学术增量 \\
    实验中KG增强检索无显著提升 & 低 & 在故障诊断这种高度结构化的知识领域，KG检索几乎必然优于纯向量检索（已有KBS 2025等多篇论文实证） \\
    课题组可用文档数据不足 & 中 & 可用公开国标/行标（GB/T系列） + 设备技术手册 + 刘老师课题组已发论文作为替代数据源 \\
    研究期间被同领域论文抢先发表 & 低 & 中文 + 动力设备（汽轮机/发电机） + 知识图谱这一交叉点极窄，竞争者很少 \\
    \bottomrule
\end{tabular}
\end{table}


% ================================================================
\section{W1--W30 逐周工作计划}
% ================================================================

每周投入 \textbf{8 小时}，分 6 阶段。整体逻辑：\textbf{搭工具（快速） → 发现问题 → 设计方案 → 对照实验 → 写论文}。

\subsection{阶段一：技术调研 + 快速搭建基础RAG（W1--W5，4月）}

\noindent\textit{阶段目标}：搞懂前沿、跑通工具链、得到一个``能用但不完美''的基础RAG——它的不完美恰好是后续研究的起点。

\noindent\textit{阶段里程碑}：{\color{tsinghua}\textbf{基础RAG效果摸底报告} → 发刘老师确认}

\begin{description}[style=nextline, leftmargin=1.2cm, labelwidth=1cm, font=\bfseries\color{tsinghua}]
\item[W1] \textbf{文献精读}。精读5篇核心论文：(1) KBS 2025 Graph RAG for Train Bogies；(2) PARAM框架（arXiv 2025）；(3) MAKG多Agent框架（2026）；(4) RAG最新综述；(5) 刘老师课题组的故障诊断代表论文。输出：2页技术调研摘要发刘老师。
\item[W2] \textbf{数据盘点与收集}。与课题组博士生沟通，梳理可用文档（技术手册、故障案例库、行业标准、论文）。整理一份``知识资源清单''，明确数据量和格式。同时下载公开国标/行标作为补充。
\item[W3] \textbf{搭建基础RAG}。用LangChain + Chroma向量数据库 + Qwen大模型（或开源替代），跑通完整Pipeline：文档加载 → 分块 → Embedding → 存储 → 检索 → 生成回答。先导入20份文档做验证。
\item[W4] \textbf{导入全部数据}。文档预处理Pipeline开发（PDF/Word → 文本 → 清洗 → 分块 → 标注元数据）。目标：处理全部可用文档（预计100--500份）。
\item[W5] \textbf{基础RAG效果摸底}。编写50道领域测试问题，跑一轮基础RAG，\textbf{手动分析失败案例}——记录``问了什么''''检索到了什么''''哪里不对''。这些失败案例就是RQ1的出发点。输出：基础RAG效果报告，发刘老师确认下一步。
\end{description}

\subsection{阶段二：构建评测数据集 + 知识图谱（W6--W10，5--6月）}

\noindent\textit{阶段目标}：建成中文QA评测数据集（论文贡献5）和故障因果知识图谱（为RQ1服务）。

\noindent\textit{阶段里程碑}：{\color{tsinghua}\textbf{200+题评测数据集 + 知识图谱初版} → 发刘老师评审}

\begin{description}[style=nextline, leftmargin=1.2cm, labelwidth=1cm, font=\bfseries\color{tsinghua}]
\item[W6] \textbf{设计评测数据集}。定义题目分类体系：按难度（简单/中等/复杂）$\times$ 类型（事实型/诊断型/方法型），每类至少30题。编写标注指南，确保每道题包含：问题、标准答案、引用来源、难度标签、类型标签。
\item[W7] \textbf{编写评测数据集（上）}。从故障案例和论文中提取100道问题。每道题人工编写标准答案和引用来源。
\item[W8] \textbf{编写评测数据集（下）}。再编写100道，完成200+题目。部分题目请课题组博士生校验。
\item[W9] \textbf{构建故障因果知识图谱}。定义本体（设备类型→部件→故障模式→原因→症状→维修方案）。用LLM从文档中半自动抽取三元组，人工校验正确性。存入Neo4j或NetworkX。
\item[W10] \textbf{阶段验证 + 汇报}。用评测数据集正式测试基础RAG，记录各项指标。把结果和知识图谱构建情况发刘老师。
\end{description}

\subsection{阶段三：核心研发——RQ1检索改进 + RQ2 Agent搭建（W11--W15，7月）}

\noindent\textit{这是最关键的阶段}：实现创新方案，初步验证有效性。

\noindent\textit{阶段里程碑}：{\color{tsinghua}\textbf{中期报告——检索改进+Agent初步对比结果} → 发刘老师}

\begin{description}[style=nextline, leftmargin=1.2cm, labelwidth=1cm, font=\bfseries\color{tsinghua}]
\item[W11] \textbf{实现混合检索}。(1) BM25关键词检索（Elasticsearch或rank\_bm25库）；(2) 知识图谱子图查询——用户Query经LLM抽取实体后查图谱关联节点；(3) 三路结果加权融合 + Reranker精排（bge-reranker模型）。
\item[W12] \textbf{实现Query改写与同义词表}。用LLM将短问题扩展为详细检索Query；构建领域同义词映射表（如``抢振''→``转子不平衡导致的强制振动''）。在评测数据集上测试改进效果。
\item[W13] \textbf{搭建多步Agent框架}。基于LangGraph实现ReAct模式Agent：Reasoning（思考）→ Action（调用工具）→ Observation（观察结果）循环。定义5个工具（见RQ2工具表）。在10道复杂题上手动调试。
\item[W14] \textbf{Agent策略设计}。设计追问逻辑（信息不足时主动求证）、置信度判断逻辑（低置信时标注``仅供参考''）、安全兜底逻辑（涉及危险操作时附加警告）。
\item[W15] \textbf{初步验证 + 向刘老师汇报}。改进RAG vs 基础RAG的初步对比；Agent vs 单轮RAG的初步对比。写中期报告。
\end{description}

\subsection{阶段四：大规模对照实验——出论文数据（W16--W20，8--9月）}

\noindent\textit{阶段目标}：用严格的对照实验回答RQ1--RQ3，产出论文所需的全部数据和图表。

\noindent\textit{阶段里程碑}：{\color{tsinghua}\textbf{三组RQ的完整实验数据 + 论文图表初稿}}

\begin{description}[style=nextline, leftmargin=1.2cm, labelwidth=1cm, font=\bfseries\color{tsinghua}]
\item[W16] \textbf{RQ1完整实验}。基础RAG vs 改进RAG，在200+题评测集上跑完整对比。记录Hit Rate@5、MRR、Chain Completeness。做消融实验：去掉KG / 去掉BM25 / 去掉Reranker，分别看效果变化。
\item[W17] \textbf{RQ2完整实验}。纯LLM vs 单轮RAG vs 多步Agent，三种系统 $\times$ 三级难度 = 9组实验。每组在对应难度的测试题上跑，记录准确率。
\item[W18] \textbf{RQ2专家评估}。邀请刘老师和2--3位博士生，对50道复杂诊断题的回答做盲评打分（1--5分：准确性、完整性、推理合理性）。计算评分者间一致性（Cohen's Kappa）。
\item[W19] \textbf{RQ3评估指标实验}。(1) 用RAGAS评分 + IFS同时评估系统；(2) 以专家评分为ground truth；(3) 计算每种指标与专家评分的Spearman相关系数。
\item[W20] \textbf{结果整理}。汇总全部实验数据，画论文图表：检索对比柱状图、Agent推理流程Case Study、IFS与RAGAS相关系数对比图、数据集统计信息表。
\end{description}

\subsection{阶段五：写论文 + 搭演示系统（W21--W25，10--11月）}

\noindent\textit{阶段里程碑}：{\color{tsinghua}\textbf{论文初稿 + 可演示的Web系统}}

\begin{description}[style=nextline, leftmargin=1.2cm, labelwidth=1cm, font=\bfseries\color{tsinghua}]
\item[W21] \textbf{搭Web演示界面}。用Streamlit或Gradio：左侧输入故障描述、右侧显示诊断回答（含引用来源链接），底部展示Agent的推理过程（检索了什么、调用了什么工具、为什么追问）。
\item[W22] 和刘老师确认论文结构和目标期刊；写第1章（动力设备智能运维背景、现有RAG在工业场景的不足、我们的研究问题）。
\item[W23] 第2章（相关工作：RAG技术、知识图谱、Agent框架在工业领域的应用现状），重点对比KBS 2025和MAKG。
\item[W24] 第3章（系统设计：基础RAG → 失败分析 → 改进方案，包含混合检索架构图、Agent工具链图、CC和IFS指标的形式化定义）。
\item[W25] 第4章（实验结果：RQ1/RQ2/RQ3的全部对比表、消融分析、Case Study、专家评估结果）。
\end{description}

\subsection{阶段六：投稿 + 结题（W26--W30，12--1月）}

\noindent\textit{阶段里程碑}：{\color{tsinghua}\textbf{论文投稿 + 鸿雁计划结题}}

\begin{description}[style=nextline, leftmargin=1.2cm, labelwidth=1cm, font=\bfseries\color{tsinghua}]
\item[W26] 第5章（结论、局限性、未来工作） + 全文润色。整理评测数据集，准备开源（GitHub + 数据说明文档）。
\item[W27] 给刘老师审改，根据反馈补充实验或修改论述。
\item[W28] 按目标期刊模板排版，提交投稿。
\item[W29] 写鸿雁计划结题报告（约15页）。
\item[W30] 提交结题材料包，项目完成。
\end{description}

\end{document}
